
\documentclass{article}
\usepackage{colortbl}
\usepackage{makecell}
\usepackage{multirow}
\usepackage{supertabular}

\begin{document}

\newcounter{utterance}

\centering \large Interaction Transcript for game `cladder', experiment `full\_v1.5\_default', episode 9213 with score{-}Qwen3.5{-}9B{-}sp.
\vspace{24pt}

{ \footnotesize  \setcounter{utterance}{1}
\setlength{\tabcolsep}{0pt}
\begin{supertabular}{c@{$\;$}|p{.15\linewidth}@{}p{.15\linewidth}p{.15\linewidth}p{.15\linewidth}p{.15\linewidth}p{.15\linewidth}}
    \# & \multicolumn{2}{c}{Player} && \multicolumn{2}{c}{Game Master} \\
    \hline

    \theutterance \stepcounter{utterance}  
    & & & \multicolumn{4}{p{0.6\linewidth}}{
        \cellcolor[rgb]{0.9,0.9,0.9}{
            \makecell[{{p{\linewidth}}}]{
                \texttt{\tiny{[P1$\langle$GM]}}
                \texttt{You are an expert in causal inference. The following question is not a typical commonsense query, but rather a meticulously designed question created by a professor specializing in causal inference, intended to assess the students' mastery of the course content.} \\
\texttt{Imagine a self{-}contained, hypothetical world with only the following conditions, and without any unmentioned factors or causal relationships: Pexu has a direct effect on rukz. Hwax has a direct effect on rukz. For those who are not pexu, the probability of rukz is 9\%. For those who are pexu, the probability of rukz is 18\%. For those who are pexu, would it be less likely to see rukz if the individual was not pexu?} \\
\texttt{Start your answer with 'Yes' or 'No', followed by additional reasoning or evidence to support your explanation.} \\
            }
        }
    }
     \\ \\

    \theutterance \stepcounter{utterance}  
    & \multicolumn{4}{p{0.6\linewidth}}{
        \cellcolor[rgb]{0.561,0.737,0.561}{
            \makecell[{{p{\linewidth}}}]{
                \texttt{\tiny{[P1$\rangle$GM]}}
                \texttt{No} \\
\\ 
\texttt{**Reasoning:**} \\
\\ 
\texttt{To answer this, we must analyze the causal structure and the specific counterfactual query within the constraints of the hypothetical world.} \\
\\ 
\texttt{1.  **Identify the Causal Structure:**} \\
\texttt{    *   The prompt states: "Pexu has a direct effect on rukz" and "Hwax has a direct effect on rukz."} \\
\texttt{    *   This implies a Directed Acyclic Graph (DAG) where both Pexu and Hwax are parents of Rukz.} \\
\texttt{    *   Crucially, the prompt specifies "without any unmentioned factors or causal relationships." This means there is no confounding variable (a common cause of both Pexu and Rukz) and no hidden bias. The association between Pexu and Rukz is purely causal.} \\
\\ 
\texttt{2.  **Analyze the Probabilities:**} \\
\texttt{    *   \$P(\textbackslash{}text\{Rukz\} | \textbackslash{}text\{Pexu\}) = 18\textbackslash{}\%\$} \\
\texttt{    *   \$P(\textbackslash{}text\{Rukz\} | \textbackslash{}neg\textbackslash{}text\{Pexu\}) = 9\textbackslash{}\%\$} \\
\texttt{    *   The probability of observing Rukz is higher for those who are Pexu (18\%) than for those who are not Pexu (9\%).} \\
\\ 
\texttt{3.  **Evaluate the Counterfactual Query:**} \\
\texttt{    *   The question asks: "For those who are pex} \\
            }
        }
    }
    & & \\ \\

    \theutterance \stepcounter{utterance}  
    & & & \multicolumn{2}{p{0.3\linewidth}}{
        \cellcolor[rgb]{0.9,0.9,0.9}{
            \makecell[{{p{\linewidth}}}]{
                \texttt{\tiny{[GM$|$GM]}}
                \texttt{No} \\
            }
        }
    }
    & & \\ \\

    \theutterance \stepcounter{utterance}  
    & & & \multicolumn{2}{p{0.3\linewidth}}{
        \cellcolor[rgb]{0.9,0.9,0.9}{
            \makecell[{{p{\linewidth}}}]{
                \texttt{\tiny{[GM$|$GM]}}
                \texttt{yes} \\
            }
        }
    }
    & & \\ \\

    \theutterance \stepcounter{utterance}  
    & & & \multicolumn{2}{p{0.3\linewidth}}{
        \cellcolor[rgb]{0.9,0.9,0.9}{
            \makecell[{{p{\linewidth}}}]{
                \texttt{\tiny{[GM$|$GM]}}
                \texttt{game\_result = LOSE} \\
            }
        }
    }
    & & \\ \\

\end{supertabular}
}

\end{document}
